\addchap{\lsAbbreviationsTitle}

We have adhered to the glossing conventions established within the Navajo corpus, in order to maintain internal consistency. Abbreviations that deviate from the Leipzig Glossing Rules are marked by an asterisk.\\

\begin{tabularx}{.45\textwidth}{lQ}
1 & first person\\
2 & second person\\
3 & third person\\
\textsc{abs*} & absolute aspect (as opposed to comparative aspect)\\
\textsc{agt} & agentive\\
\textsc{anim} & animate person\\
\textsc{ams} & aspect, mood and subject\\
\textsc{ano} & (single) animate object\\
\textsc{ar} & areal\\
\textsc{bi} & bi-object (occurs only as \textsc{3}.\textsc{bi})\\
\textsc{comp*} & comparative aspect\\
\textsc{con} & conative aspect\\
\textsc{conc} & conclusive aspect\\
\textsc{cont} & continuative aspect\\
\textsc{curs} & cursive aspect\\
\textsc{d} & d-classifier (includes d-segment of first person duoplural forms)\\
\textsc{dir} & directive\\
\textsc{distr} & distributive\\
\textsc{du} & dual\\
\textsc{du}/\textsc{pl}&duoplural (dual or plural)\\
\textsc{dur} & durative\\
\textsc{emph} & emphatic\\
\end{tabularx}
\begin{tabularx}{.50\textwidth}{lQ}
\textsc{ffo} & flat flexible object\\
\textsc{foc} & focus\\
\textsc{fut} & future\\
\textsc{gen*} & generic\\
\textsc{imp} & imperative\\
\textsc{impf*} & imperfective\\
\textsc{inc} & inceptive\\
\textsc{inch} & inchoative\\
\textsc{indef} & indefinite\\
\textsc{intrg} & interrogative\\
\textsc{ints} & intensifier\\
\textsc{iter} & iterative\\
\textsc{l} & \textit{l}-classifier (incl. voicing of \textit{ł}-classifier)\\
\textsc{ł} & \textit{ł}-classifier (incl. devoicing of \textit{l}-classifier)\\
\textsc{lpb} & load, pack, burden\\
\textsc{mom}1 & momentaneous aspect (1)\\
\textsc{mom}2 & momentaneous aspect (2)\\
\textsc{mom}3 & momentaneous aspect (3)\\
{n*} & noun or nominal root\\
\textsc{neu} & neuter aspect\\
\textsc{ni} & \textit{ni}-aspectival (specifies \textsc{perf} or \textsc{impf} aspect)\\
\textsc{nmlz} & nominalizer\\
\textsc{nsp} & non-specific\\
\textsc{obj} & object\\
\textsc{oc} & open container\\
\textsc{opt} & optative\\
\end{tabularx}

\newpage
\begin{tabularx}{.45\textwidth}{lQ}
\textsc{osbj} & outer subject\\
\textsc{pas*} & passive\\
\textsc{pcl} & particular\\
\textsc{poss} & possessive\\
\textsc{pcl} & particular\\
\textsc{perf*} & perfective\\
\textsc{pl} & plural\\
{pro*} & pronoun or pronominal root\\
{post*} & postposition or postpositional root\\
\textsc{pres*} & present\\
\textsc{prog} & progressive\\
\textsc{prvb} & preverb \\
\textsc{pst} & past\\
\textsc{qual} & qualifier\\
\textsc{rec*} & reciprocal\\
\textsc{ref} & referential\\
\textsc{refl} & reflexive\\
\textsc{rep} & repetitive\\
\textsc{rev} & reversionary\\
\textsc{rec} & reciprocal\\
\textsc{refl} & reflexive\\
\textsc{rep} & repetitive\\
\textsc{rev} & reversionary\\
\\
\end{tabularx}
\begin{tabularx}{.45\textwidth}{lQ}
\textsc{sbj} & subject\\
\textsc{semf} & semelfactive\\
\textsc{ser} & seriative\\
\textsc{sfo} & slender flexible object\\
\textsc{sg} & singular\\
\textsc{si} & \textit{si}-aspectival (specifies \textsc{perf} or \textsc{impf} aspect)\\
\textsc{sitr} & semeliterative\\
\textsc{sp*} & specific\\
\textsc{sro} & solid roundish object\\
\textsc{sso} & stiff slender object\\
\textsc{sub} & subordination\\
\textsc{term} & terminative\\
\textsc{th} & thematic (unclear lexical semantics)\\
\textsc{trnsl} & transitional\\
\textsc{usit} & usitative aspect\\
{v*} & verb or verbal root\\
\textsc{val} & valence\\
\textsc{yi}\textsubscript{1} & \textit{yi}-object (occurs only as \textsc{3}.\textsc{yi})\\
\textsc{yi}\textsubscript{2} & \textit{yi}-aspectival (specifies \textsc{perf} or \textsc{impf} aspect)\\
\textsc{zero} & \textit{zero}-aspectival (specifies \textsc{perf} or \textsc{impf} aspect)\\
\end{tabularx}
